\chapter{T\'itulo do primeiro cap\'itulo} \label{cap:cap1}

Este \'e o primeiro cap\'itulo. 

\section{Introdu\c{c}\~ao}\label{sec:intro}

Esta seção deve ser uma extensão do resumo. Aqui o aluno deve ter um primeiro parágrafo falando sobre o contexto, isto irá situar o leitor, ele irá saber logo de início o assunto abordado no trabalho.

Após posicionar o leitor, o segundo paragrafo deve levantar o contexto problemático, ou seja, os problemas que ele identificou ou a literatura corrente relata (outros trabalhos).

O terceiro parágrafo deve ter um gancho com o segundo. Dado o contexto problemático, este trabalho teve como objetivo (por claramente o objetivo). Aqui o aluno deve deixar bem claro qual o objetivo do trabalho. Com isso, o leitor irá saber exatamente o que o trabalho irá apresentar no final, já que o objetivo está totalmente ligado aos resultados. Neste terceiro parágrafo, o aluno ainda pode dizer como atingiu este objetivo, qual método de pesquisa ele utilizou, como ele realizou a pesquisa, de maneira sucinta, tendo em vista que isso será relatado em detalhes em uma seção posterior.

O quarto parágrafo deve relatar brevemente os principais achados do trabalho. Deve conter um pequeno resumo das principais contribuições da pesquisa. Caso tenha sido uma pesquisa com 50 alunos, aqui deve ter algo como: realizou-se uma pesquisa com 50 alunos, foi verificado (veja que a escrita é sempre de maneira impessoal) que 80\textdiscount  destes tiveram um bom rendimento nas atividades repassadas, obtendo X pontos, etc. Lembrando que isso são exemplos, cada caso terá um resultado específico. 

O último parágrafo deve ter a organização do artigo. Ex. Além desta seção de introdução, este artigo está organizado em X seções. A Seção 2 apresenta os conceitos necessários para o entendimento do trabalho. A Seção 3 mostra o método de pesquisa que foi seguido para atingir os objetivos. Seção 4...


Exemplo de Figura: Ver Figura~\ref{fig:exefig}.

\begin{figure}[!ht]
\centering
\includegraphics[width=0.3\linewidth]{figuras/exefig.eps}
\caption{Exemplo de figura}
\label{fig:exefig}
\end{figure}

Este \'e o primeiro cap\'itulo. 

Exemplo de refer\^encias~\cite{tese2011, confinter2011, confnac2011, rev2011, site2011}.

Este \'e o primeiro cap\'itulo. 

Exemplo de Tabela: ver Tabela~\ref{tab:exetab}.

\begin{table}[!ht]
\begin{center}
\caption{Distribuição IMC em Adultos}
\label{tab:exetab}
\begin{tabular}{|c |c |}
\hline
\textbf{\textbf{Classificação}} & \textbf{IMC}\\
\hline\hline
Baixo Peso & < 18,5 \\
Peso Adequado & > 18,5 e < 25 \\
Sobrepeso & > 25 e < 30 \\
Obesidade & > 30 \\
\hline
\end{tabular}
\end{center}
\end{table}

Este \'e o primeiro cap\'itulo. 
 


\section{Desenvolvimento}\label{sec:desenv}

Conforme explicado na Se\c{c}\~ao~\ref{sec:intro}, nesta segunda se\c{c}\~ao ser\'a descrito o desenvolvimento.
